\section*{अध्याय \ref{ch:g9:elements-universe-earth} --- ब्रह्मांड में और पृथ्वी पर तत्व}

\begin{solution}{exo:g9:elements-universe-earth:1}
हाइड्रोजन (लगभग \qty{74}{\percent}) और हीलियम (लगभग
\qty{25}{\percent})।
\end{solution}

\begin{solution}{exo:g9:elements-universe-earth:2}
दोनों में ऑक्सीजन।
\end{solution}

\begin{solution}{exo:g9:elements-universe-earth:3}
तारों के भीतर, और उनके विस्फोटों में।
\end{solution}

\begin{solution}{exo:g9:elements-universe-earth:4}
लगभग \qty{3.5}{\percent}।
\end{solution}

\begin{solution}{exo:g9:elements-universe-earth:5}
$46.6 + 27.7 = \qty{74.3}{\percent}$: लगभग तीन-चौथाई।
\end{solution}

\begin{solution}{exo:g9:elements-universe-earth:6}
वे बहुत हल्की गैसें हैं: वे नवजात पृथ्वी से अंतरिक्ष में निकल भागीं
(और हीलियम कोई यौगिक नहीं बनाती जो उसे चट्टानों में रोक पाता)।
\end{solution}

\begin{solution}{exo:g9:elements-universe-earth:7}
$0.23 \times 70 = \qty{16.1}{kg}$।
\end{solution}

\begin{solution}{exo:g9:elements-universe-earth:8}
ऑक्सीजन \omterm{def:g7:atoms-and-molecules:atom}{परमाणु} हाइड्रोजन \omterm{def:g7:atoms-and-molecules:atom}{परमाणुओं} से 16 गुना भारी हैं: कम ऑक्सीजन \omterm{def:g7:atoms-and-molecules:atom}{परमाणु} भी
भार में अधिक होते हैं।
\end{solution}

\begin{solution}{exo:g9:elements-universe-earth:9}
लगभग $0.081 \times 1000 = \qty{81}{kg}$ ऐलुमिनियम।
\end{solution}

\begin{solution}{exo:g9:elements-universe-earth:10}
ऑक्सीजन और कैल्सियम (और, भूपर्पटी के आठ मुख्य तत्वों के अलावा,
हाइड्रोजन, कार्बन और फ़ॉस्फ़ोरस भी भूपर्पटी में थोड़ी मात्रा में
मौजूद हैं)।
\end{solution}

\begin{solution}{exo:g9:elements-universe-earth:11}
$40 \times 1.67 \times 10^{-27} = \qty{6.68e-26}{kg}$;
$1.06 / 6.68 \times 10^{-26} \approx 1.6 \times 10^{25}$ कैल्सियम \omterm{def:g7:atoms-and-molecules:atom}{परमाणु}।
\end{solution}

\begin{solution}{exo:g9:elements-universe-earth:12}
\omterm{def:g5:irreversible-changes:chemical-change}{रासायनिक परिवर्तनों} में \omterm{def:g7:atoms-and-molecules:atom}{परमाणु} न बनते हैं, न नष्ट होते हैं: वे केवल एक यौगिक से
दूसरे में जाते हैं (चट्टान, गैस, शर्करा, सजीव शरीर)। पृथ्वी के
कार्बन \omterm{def:g7:atoms-and-molecules:atom}{परमाणु} बार-बार काम आते हैं।
\end{solution}

\begin{solution}{pb:g9:elements-universe-earth:1}
\textbf{1.} $0.61 \times 70 = \qty{42.7}{kg}$।

\textbf{2.} कार्बन $0.23 \times 70 = \qty{16.1}{kg}$; हाइड्रोजन
$0.10 \times 70 = \qty{7.0}{kg}$।

\textbf{3.} $61 + 23 + 10 = \qty{94}{\percent}$।

\textbf{4.} पानी, \ce{H2O}।

\textbf{5.} $16 \times 1.67 \times 10^{-27} = \qty{2.67e-26}{kg}$।

\textbf{6.} कार्बन: $12 \times 1.67 \times 10^{-27} = \qty{2.00e-26}{kg}$;
हाइड्रोजन: \qty{1.67e-27}{kg}।

\textbf{7.} 16।

\textbf{8.} \omterm{def:g9:inside-the-atom:nucleus}{इलेक्ट्रॉन} \omterm{def:g9:inside-the-atom:proton}{न्यूक्लिऑन} से लगभग 1836 गुना हल्का है: \omterm{def:g9:inside-the-atom:nucleus}{इलेक्ट्रॉन}
द्रव्यमान के हज़ारवें भाग से भी कम हैं।

\textbf{9.} $7.0 / 1.67 \times 10^{-27} \approx 4.2 \times 10^{27}$
हाइड्रोजन \omterm{def:g7:atoms-and-molecules:atom}{परमाणु}।

\textbf{10.} $16.1 / 2.00 \times 10^{-26} \approx 8.0 \times 10^{26}$
कार्बन \omterm{def:g7:atoms-and-molecules:atom}{परमाणु}।

\textbf{11.} $42.7 / 2.67 \times 10^{-26} \approx 1.60 \times 10^{27}$
ऑक्सीजन \omterm{def:g7:atoms-and-molecules:atom}{परमाणु}।

\textbf{12.} हमारी गिनती, $1.60 \times 10^{27}$, अनुमान
$1.61 \times 10^{27}$ से मेल खाती है: \qty{70}{kg} के शरीर में लगभग $1.6 \times
10^{27}$ ऑक्सीजन \omterm{def:g7:atoms-and-molecules:atom}{परमाणु} होते हैं।
\end{solution}
